%=====================================================================
\chapter{Grouping, Aggregation and Window Functions}\label{ch:grouping}
%=====================================================================
\locator{4}{Part IV --- Asking Questions}

\begin{outcomes}
By the end of this chapter you will be able to:
\begin{tight}
  \item \textbf{Use} \sqlkw{COUNT}, \sqlkw{SUM}, \sqlkw{AVG}, \sqlkw{MIN} and
        \sqlkw{MAX}, and \textbf{predict} what each does with a null.
  \item \textbf{Write} a \sqlkw{GROUP BY} query and \textbf{state} which
        columns may appear in its \sqlkw{SELECT} list.
  \item \textbf{Choose} between \sqlkw{WHERE} and \sqlkw{HAVING} by what the
        condition is about.
  \item \textbf{Write} a window function with \sqlkw{PARTITION BY} and
        \sqlkw{ORDER BY}, and \textbf{explain} how it differs from grouping.
  \item \textbf{Read} a query plan for \sqlkw{TEMP B-TREE}, and \textbf{say}
        what index would remove it.
\end{tight}
\end{outcomes}

\begin{prereq}
\chref{select} for three-valued logic, which the aggregates handle in their
own way. \chref{subqueries} for correlated sub-queries, one of which this
chapter replaces. You need \texttt{code/m16\_grouping.py}.
\end{prereq}

\begin{keyterms}
aggregate function $\bullet$ \sqlkw{COUNT} $\bullet$ \sqlkw{SUM} $\bullet$
\sqlkw{AVG} $\bullet$ \sqlkw{MIN} $\bullet$ \sqlkw{MAX} $\bullet$
\sqlkw{GROUP BY} $\bullet$ group $\bullet$ \sqlkw{HAVING} $\bullet$ bare
column $\bullet$ window function $\bullet$ \sqlkw{OVER} $\bullet$
\sqlkw{PARTITION BY} $\bullet$ \sqlkw{RANK} $\bullet$ \sqlkw{DENSE\_RANK}
$\bullet$ \sqlkw{ROW\_NUMBER} $\bullet$ running total $\bullet$
\sqlkw{TEMP B-TREE}
\end{keyterms}

\section{Collapsing Rows Into One}

Every query so far has returned rows that existed --- possibly filtered,
possibly joined, but each one traceable to a tuple in the database. An
\term{aggregate function} returns a row that existed nowhere: the \emph{average} grade, the
\emph{number} of students.

\begin{definitionbox}[The five aggregates, and what they skip]
\begin{tabular}{@{}L{3.6cm}L{10.4cm}@{}}
\sqlkw{COUNT(*)} & the number of rows. Counts everything, nulls included. \\
\sqlkw{COUNT}$(x)$ & the number of rows where $x$ is \textbf{not null}. \\
\sqlkw{SUM}, \sqlkw{AVG} & over the non-null values only. \sqlkw{AVG}
  divides by the number of non-null values, not by the number of rows. \\
\sqlkw{MIN}, \sqlkw{MAX} & over the non-null values only. \\
\end{tabular}

\smallskip
Every aggregate except \sqlkw{COUNT(*)} \textbf{skips nulls}. Over a set that
is entirely null, \sqlkw{SUM} and \sqlkw{AVG} return null --- not zero ---
and \sqlkw{COUNT} returns 0.
\end{definitionbox}

\section{\sqlkw{GROUP BY}}

\begin{definitionbox}[Grouping]
\sqlkw{GROUP BY} $a$ partitions the rows into groups, one for each distinct
value of $a$, and an aggregate in the \sqlkw{SELECT} list is then computed
\emph{per group}. The result has one row per group.

\smallskip
\textbf{The rule for the \sqlkw{SELECT} list:} every column in it must either
be in the \sqlkw{GROUP BY}, or be inside an aggregate. A column that is neither
--- a \term{bare column} --- has no single value per group, so there is
nothing correct to print.

\smallskip
Null forms a group of its own. All the nulls are grouped together, even though
\texttt{NULL = NULL} is \unk{}: grouping uses ``is the same value''
rather than ``compares equal''.
\end{definitionbox}

\begin{alertbox}[SQLite allows the bare column, and it should not]
\begin{lstlisting}[language=sql]
SELECT programme, name, COUNT(*) FROM student GROUP BY programme;
\end{lstlisting}

PostgreSQL rejects this: \texttt{name} is neither grouped nor aggregated.
SQLite runs it and prints, for each programme, the name of \emph{some} student
in it --- whichever row the engine happened to process last. The output looks
like a meaningful answer, and on each run it may name a different student.

\smallskip
There is one documented exception: with exactly one \sqlkw{MIN} or
\sqlkw{MAX}, SQLite guarantees the bare columns come from the row holding that
minimum or maximum. It is a convenience no other system offers, so code that
depends on it will not move. \textbf{In this course, every non-aggregated
column is in the \sqlkw{GROUP BY}.}
\end{alertbox}

\section{Measured: Nulls, Plans, and a Rank}

\begin{measured}[What aggregates skip, what an index removes, and what a window replaces]
\textsf{\footnotesize \texttt{code/m16\_grouping.py}, seeded database,
SQLite 3.51.2, Intel i5-6600, Windows 11. Best of three runs.}

\medskip
\textbf{1. Aggregates and nulls} --- \textit{enrolment}, 60{,}037 rows

\begin{tabular}{@{}L{6.4cm}R{2.4cm}L{5.4cm}@{}}
\sqlkw{COUNT(*)}      & 60{,}037 & every row \\
\sqlkw{COUNT}(grade)  & 40{,}045 & rows where grade is not null \\
difference            & 19{,}992 & the ungraded FALL2026 rows \\
\sqlkw{AVG}(points)   & \textbf{2.383} & nulls skipped: $\div$ 40{,}045 \\
\sqlkw{AVG}(\sqlkw{COALESCE}(points, 0)) & \textbf{1.590} & nulls as zero:
  $\div$ 60{,}037 \\
\end{tabular}

\smallskip
The second average is \textbf{33\% lower}, and it has decided that an ungraded
course is an F. \sqlkw{GROUP BY} grade puts the 19{,}992 nulls in a group of
their own, beside A (7{,}268), B (12{,}756), C (11{,}283), D (5{,}524) and F
(3{,}214).

\medskip
\textbf{2. \sqlkw{GROUP BY} with and without a supporting index} --- enrolments
per course

\begin{tabular}{@{}R{2.0cm}R{2.4cm}R{2.6cm}R{2.6cm}R{2.4cm}@{}}
\textbf{students} & \textbf{enrolments} & \textbf{no index} & \textbf{with index}
& \textbf{speed-up} \\
4{,}000  & 60{,}037  & 37.0\,ms  & 6.8\,ms  & 5.5$\times$ \\
8{,}000  & 120{,}134 & 86.3\,ms  & 13.9\,ms & 6.2$\times$ \\
16{,}000 & 240{,}025 & 174.2\,ms & 31.3\,ms & 5.6$\times$ \\
\end{tabular}

\smallskip
Without the index the plan says \texttt{USE TEMP B-TREE FOR GROUP BY}: the
rows were sorted into groups before they could be counted. An index on
\texttt{course\_code} is already in that order, so the sort disappears from the
plan --- and five-sixths of the cost goes with it. The speed-up is roughly
constant as the table grows, because both forms are linear and the sort was a
fixed proportion of the work.

\medskip
\textbf{3. The same filter in \sqlkw{WHERE} and in \sqlkw{HAVING}}

\begin{tabular}{@{}L{6.0cm}R{1.5cm}R{1.7cm}L{3.0cm}@{}}
& \textbf{ms} & \textbf{groups} & \\
\texttt{WHERE course\_code LIKE 'CS-\%'}  & 13.1 & 6 & \\
\texttt{HAVING course\_code LIKE 'CS-\%'} & 13.9 & 6 & same answer \\
\end{tabular}

\smallskip
\textbf{The two compiled to the identical program}, 47 bytecode operations
each, with the \sqlkw{LIKE} tested at operation 8, \emph{before} the grouping
sort at operation 12. SQLite moved the \sqlkw{HAVING} condition into
\sqlkw{WHERE} by itself --- which it can do only because the condition mentions
no aggregate. \texttt{HAVING COUNT(*) > 1000} was tested at operation 33,
\emph{after} the sort at operation 9, because a count does not exist until the
grouping is done.

\medskip
\textbf{4. A rank within each programme} --- by number of A grades

\begin{tabular}{@{}R{2.0cm}R{3.4cm}R{3.4cm}R{2.0cm}L{2.4cm}@{}}
\textbf{students} & \textbf{window \sqlkw{RANK()}} & \textbf{correlated}
& \textbf{ratio} & \textbf{agree?} \\
1{,}000 & 16.9\,ms  & 51.6\,ms  & 3$\times$ & yes \\
2{,}000 & 38.4\,ms  & 190.4\,ms & 5$\times$ & yes \\
4{,}000 & 103.2\,ms & 736.7\,ms & \textbf{7$\times$} & yes \\
\end{tabular}

\smallskip
The ratio grows with the data. The correlated form counts, for each student,
how many in the same programme did better --- every student against every
other in the programme, quadratic. The window function sorts each programme
once and reads the rank off the order.
\end{measured}

\begin{alertbox}[What the four results say about how to write these queries]
\begin{tight}
  \item \textbf{Never \sqlkw{COALESCE} to zero inside an aggregate unless a
        missing value really is a zero.} It moved the average by a third.
        That is the same fact \chref{ddl} found with a text value in a numeric
        column, approached from the other side.
  \item \textbf{\sqlkw{TEMP B-TREE} in a plan is a sort you could have avoided.}
        An index on the grouping column supplies the order for free.
        \chref{indexes} shows where the same idea stops working.
  \item \textbf{Put a condition on a column in \sqlkw{WHERE} because that is
        what it means, not because it is faster here.} SQLite fixed the
        \sqlkw{HAVING} version itself; that is an optimisation, not a promise,
        and the next system you use may not make it.
  \item \textbf{For rankings, running totals and ``top $n$ per group'', use a
        window function.} The correlated alternative was seven times slower at
        4{,}000 students and getting worse.
\end{tight}
\end{alertbox}

\section{Window Functions}

\begin{definitionbox}[A window function]
A \term{window function} computes a value for each row from a set of related
rows --- its \emph{window} --- \textbf{without collapsing them}. Grouping turns
many rows into one; a window function keeps every row and adds a column.

\begin{lstlisting}[language=sql]
SELECT roll_no, programme, n_a,
       RANK() OVER (PARTITION BY programme ORDER BY n_a DESC) AS pos
  FROM per_student;
\end{lstlisting}

\smallskip
\sqlkw{PARTITION BY} divides the rows into windows, like \sqlkw{GROUP BY} but
without merging them. \sqlkw{ORDER BY} inside \sqlkw{OVER} orders each window.
The function then runs along the ordered window.
\end{definitionbox}

\dsfig{%
\begin{tikzpicture}[
  cell/.style={draw=primary!35, line width=0.5pt, minimum width=1.45cm,
               minimum height=0.40cm, font=\scriptsize\ttfamily,
               text=primary!65!black, inner sep=0pt},
  hd/.style={cell, fill=primary, text=white, font=\scriptsize\bfseries},
  p1/.style={cell, fill=secondary!14},
  p2/.style={cell, fill=goldhl!16},
  nw/.style={cell, fill=greenhl!20, font=\scriptsize\bfseries\ttfamily,
             text=greenhl!55!black},
  ttl/.style={font=\scriptsize\bfseries, text=#1, anchor=south},
  arr/.style={-{Stealth[length=2.2mm]}, primary!50, line width=0.9pt},
  nt/.style={font=\scriptsize\itshape, text=neutral, align=center}
]
  % source rows
  \node[ttl=primary] at (-5.3,0.5) {per\_student};
  \foreach \c [count=\i] in {prog, n\_a}{\node[hd] at (-6.75+\i*1.45,0.2) {\c};}
  \foreach \r/\p/\n/\s in {1/BSIT/9/p1, 2/BSIT/7/p1, 3/BSIT/5/p1,
                           4/BSCS/8/p2, 5/BSCS/6/p2}{
    \node[\s] at (-5.30,-0.4*\r+0.2) {\p};
    \node[\s] at (-3.85,-0.4*\r+0.2) {\n};}

  % group by
  \draw[arr] (-2.95,-0.25) -- (-1.75,0.15);
  \node[ttl=secondary!70!black] at (0.0,0.5) {\sqlkw{GROUP BY} prog};
  \foreach \c [count=\i] in {prog, MAX(n\_a)}{\node[hd] at (-1.45+\i*1.45,0.2) {\c};}
  \node[p1] at (0.0,-0.2) {BSIT};  \node[p1] at (1.45,-0.2) {9};
  \node[p2] at (0.0,-0.6) {BSCS};  \node[p2] at (1.45,-0.6) {8};
  \node[nt, text width=3.6cm] at (0.72,-1.35)
       {five rows become \textbf{two}: one per group};

  % window
  \draw[arr] (-2.95,-1.15) .. controls (-1.2,-2.4) and (2.5,-2.4) .. (3.4,-1.3);
  \node[ttl=greenhl!60!black] at (5.4,0.5)
       {\sqlkw{RANK() OVER} (\sqlkw{PARTITION BY} prog)};
  \foreach \c [count=\i] in {prog, n\_a, rank}{\node[hd] at (2.95+\i*1.45,0.2) {\c};}
  \foreach \r/\p/\n/\k/\s in {1/BSIT/9/1/p1, 2/BSIT/7/2/p1, 3/BSIT/5/3/p1,
                              4/BSCS/8/1/p2, 5/BSCS/6/2/p2}{
    \node[\s] at (4.40,-0.4*\r+0.2) {\p};
    \node[\s] at (5.85,-0.4*\r+0.2) {\n};
    \node[nw] at (7.30,-0.4*\r+0.2) {\k};}
  \node[nt, text width=4.4cm] at (5.85,-2.35)
       {five rows stay \textbf{five}; each gains a column computed from its
        own partition};
\end{tikzpicture}}%
{The difference in one picture. Grouping collapses each partition into a single
row; a window function keeps every row and adds a value computed from the
partition it belongs to. The colours mark the two partitions.}{groupwindow}

\begin{examplebox}[\sqlkw{ROW\_NUMBER}, \sqlkw{RANK} and \sqlkw{DENSE\_RANK} on a tie]
Four students in one programme with 9, 7, 7 and 5 A grades:

\medskip
\begin{tabular}{@{}R{2.2cm}R{3.2cm}R{2.6cm}R{3.2cm}@{}}
\textbf{A grades} & \sqlkw{ROW\_NUMBER} & \sqlkw{RANK} & \sqlkw{DENSE\_RANK} \\
9 & 1 & 1 & 1 \\
7 & 2 & 2 & 2 \\
7 & 3 & 2 & 2 \\
5 & 4 & \textbf{4} & \textbf{3} \\
\end{tabular}

\smallskip
\sqlkw{ROW\_NUMBER} never ties, so it numbers the two sevens arbitrarily ---
use it only when the order is fully determined. \sqlkw{RANK} gives ties the
same number and \emph{skips} the next, as a sports table does.
\sqlkw{DENSE\_RANK} does not skip. The university's merit list is a policy
decision between the last two, and the query has to make it explicitly.
\end{examplebox}

\begin{worked}[A running total of credit hours, per student, in term order]
\textbf{The question.} For each enrolment of one student, how many credit hours
had they completed up to and including that term?

\medskip
\textbf{Step 1 --- why grouping cannot do it.} \sqlkw{GROUP BY} roll\_no gives
one total per student. The question wants every row kept, each with its own
total so far.

\medskip
\textbf{Step 2 --- the window.}
\begin{lstlisting}[language=sql]
SELECT e.semester, c.course_code, c.credit_hours,
       SUM(c.credit_hours) OVER (
         PARTITION BY e.roll_no
         ORDER BY e.semester
         ROWS BETWEEN UNBOUNDED PRECEDING AND CURRENT ROW
       ) AS hours_so_far
  FROM enrolment e JOIN course c ON c.course_code = e.course_code
 WHERE e.roll_no = (SELECT MIN(roll_no) FROM student);
\end{lstlisting}

\medskip
\textbf{Step 3 --- the frame.} \sqlkw{ROWS BETWEEN UNBOUNDED PRECEDING AND
CURRENT ROW} says ``from the first row of the window to this one''. That is
what makes it a running total rather than a grand total repeated on every row.

\medskip
\textbf{Step 4 --- a trap in the ordering.} \texttt{ORDER BY e.semester} sorts
the text \texttt{'FALL2025'}, \texttt{'FALL2026'}, \texttt{'SPRING2026'}
\emph{alphabetically} --- and so puts Fall 2026 before Spring 2026, which is
wrong. A term label that sorts correctly as text has to be designed to,
\texttt{'2025-2'} rather than \texttt{'FALL2025'}. This is the same column that
broke \sqlkw{NATURAL JOIN} in \chref{joins}, and the second reason to revisit
it when \chref{normalforms} reconsiders the schema.
\end{worked}

\begin{pitfall}
\begin{tight}
  \item \textbf{A bare column in a grouped query.} SQLite prints some row's
        value; other systems refuse.
  \item \textbf{\sqlkw{COALESCE} to zero inside \sqlkw{AVG}.} Measured: a third
        off.
  \item \textbf{Expecting \sqlkw{SUM} of nothing to be 0.} It is null.
        \sqlkw{COALESCE}(\sqlkw{SUM}(x), 0) is the honest fix there.
  \item \textbf{\sqlkw{HAVING} for a condition on a column.} It works here,
        because SQLite rewrites it. Write \sqlkw{WHERE}.
  \item \textbf{A correlated sub-query for a rank.} Measured: seven times
        slower and growing. Use \sqlkw{RANK} \sqlkw{OVER}.
  \item \textbf{\sqlkw{ROW\_NUMBER} where ties are possible.} It breaks them
        arbitrarily.
  \item \textbf{Ordering by a term label stored as text.} It sorts
        alphabetically, which is not chronologically.
\end{tight}
\end{pitfall}

\begin{lab}[Move an average by a third, remove a sort, and rank a programme]
\begin{lstlisting}[language=bash]
# --- 1. The seeded database ----------------------------------------
python code/schema/seed.py university.db
sqlite3 university.db
\end{lstlisting}
\begin{lstlisting}[language=sql]
-- --- 2. Count three ways ----------------------------------------
.headers on
.mode box
.timer on
SELECT COUNT(*), COUNT(grade), COUNT(DISTINCT grade) FROM enrolment;
-- Three numbers. Say what each one counts.

-- --- 3. Two averages, one of them a decision about failure ------
SELECT AVG(CASE grade WHEN 'A' THEN 4 WHEN 'B' THEN 3 WHEN 'C' THEN 2
                      WHEN 'D' THEN 1 WHEN 'F' THEN 0 END) FROM enrolment;
SELECT AVG(COALESCE(CASE grade WHEN 'A' THEN 4 WHEN 'B' THEN 3
              WHEN 'C' THEN 2 WHEN 'D' THEN 1 WHEN 'F' THEN 0 END, 0))
  FROM enrolment;
-- By what percentage do they differ, and what did the second decide?

-- --- 4. The null group -----------------------------------------
SELECT grade, COUNT(*) FROM enrolment GROUP BY grade ORDER BY grade;

-- --- 5. A sort you can see --------------------------------------
EXPLAIN QUERY PLAN
SELECT course_code, COUNT(*) FROM enrolment GROUP BY course_code;
SELECT course_code, COUNT(*) FROM enrolment GROUP BY course_code;
-- Find TEMP B-TREE. Record the time.

-- --- 6. Remove it -----------------------------------------------
CREATE INDEX idx_enr_course ON enrolment(course_code);
EXPLAIN QUERY PLAN
SELECT course_code, COUNT(*) FROM enrolment GROUP BY course_code;
SELECT course_code, COUNT(*) FROM enrolment GROUP BY course_code;
-- What disappeared from the plan, and what happened to the time?

-- --- 7. WHERE against HAVING, and a condition that cannot move ---
SELECT course_code, COUNT(*) FROM enrolment
 GROUP BY course_code HAVING COUNT(*) > 1300;
-- Try to write that one with WHERE instead. Why can you not?

-- --- 8. A rank, with a window ------------------------------------
WITH a AS (SELECT s.roll_no, s.programme, SUM(e.grade = 'A') AS n_a
             FROM student s JOIN enrolment e ON e.roll_no = s.roll_no
            GROUP BY s.roll_no)
SELECT programme, roll_no, n_a,
       RANK()       OVER w AS rnk,
       DENSE_RANK() OVER w AS dense
  FROM a WINDOW w AS (PARTITION BY programme ORDER BY n_a DESC)
 ORDER BY programme, rnk LIMIT 12;
-- Find a tie. Show where RANK and DENSE_RANK disagree.
\end{lstlisting}
\begin{lstlisting}[language=bash]
# --- 9. The measured version ----------------------------------------
python code/m16_grouping.py
\end{lstlisting}

\smallskip
\textbf{What to notice.} Step 7 is the real difference between the two
clauses. A condition on \texttt{course\_code} could go in either, and SQLite
quietly makes them identical. A condition on \sqlkw{COUNT(*)} can only go in
\sqlkw{HAVING}, because when \sqlkw{WHERE} runs the count has not been
computed. Choose by \emph{what the condition is about}, and the clause chooses
itself.
\end{lab}

\begin{checkpoint}
\begin{tightnum}
  \item On the seeded data \sqlkw{AVG}(points) is 2.383 and
        \sqlkw{AVG}(\sqlkw{COALESCE}(points, 0)) is 1.590. Explain both
        numbers in terms of what each one divides by, and say which one the
        examinations office should publish and why.
  \item The \sqlkw{WHERE} and \sqlkw{HAVING} versions of the same filter
        compiled to identical programs. Explain why SQLite was able to do that
        for \texttt{course\_code LIKE 'CS-\%'} and not for \texttt{COUNT(*) >
        1000}, and say why you should write \sqlkw{WHERE} anyway.
  \item Give a question that needs a window function rather than
        \sqlkw{GROUP BY}, write the query, and say what \sqlkw{GROUP BY} would
        have lost.
\end{tightnum}
\end{checkpoint}

\begin{chaptersummary}
\begin{tight}
  \item Every aggregate except \sqlkw{COUNT(*)} \textbf{skips nulls};
        \sqlkw{AVG} divides by the non-null count; \sqlkw{SUM} of nothing is
        null, not 0.
  \item \textbf{Measured:} \sqlkw{AVG} skipping nulls gave 2.383;
        \sqlkw{COALESCE} to zero gave 1.590 --- \textbf{a third lower}, and a
        decision that ungraded means failed.
  \item \sqlkw{GROUP BY} makes one row per group. Every \sqlkw{SELECT} column
        must be grouped or aggregated; SQLite permits a bare column and prints
        an arbitrary row's value.
  \item Null forms its own group.
  \item \textbf{Measured:} an index on the grouping column removed
        \sqlkw{TEMP B-TREE FOR GROUP BY} from the plan and made the query
        \textbf{5.5--6.2$\times$} faster.
  \item \textbf{Measured:} the \sqlkw{WHERE} and \sqlkw{HAVING} forms of a
        column condition \textbf{compiled to the same program} --- SQLite
        moved it. An aggregate condition cannot move. Choose by what the
        condition is about.
  \item A \textbf{window function} adds a value to every row without
        collapsing them. \sqlkw{PARTITION BY} defines the windows,
        \sqlkw{ORDER BY} orders them, the frame bounds them.
  \item \textbf{Measured:} a window \sqlkw{RANK()} was \textbf{7$\times$}
        faster than the correlated equivalent at 4{,}000 students, and the gap
        was widening.
  \item \sqlkw{RANK} skips after a tie, \sqlkw{DENSE\_RANK} does not,
        \sqlkw{ROW\_NUMBER} breaks ties arbitrarily.
  \item A term label stored as \texttt{'FALL2025'} sorts alphabetically, not
        chronologically --- the second fault found in the \texttt{semester}
        column.
\end{tight}
\end{chaptersummary}

\begin{reviewq}
  \item \textbf{(MCQ)} On a column with 10 rows of which 4 are null,
        \sqlkw{COUNT}$(x)$ returns: (a)~10; (b)~6; (c)~4; (d)~null.
  \item \textbf{(MCQ)} A condition on \sqlkw{COUNT(*)} must go in:
        (a)~\sqlkw{WHERE}; (b)~\sqlkw{HAVING}; (c)~\sqlkw{ORDER BY};
        (d)~either.
  \item \textbf{(Short)} Why can a bare column in a grouped query not have a
        correct value, and what does SQLite print instead?
  \item \textbf{(Short)} What does \sqlkw{TEMP B-TREE FOR GROUP BY} in a plan
        mean, and what removes it?
  \item \textbf{(Short)} Give four rows with a tie and show how
        \sqlkw{RANK}, \sqlkw{DENSE\_RANK} and \sqlkw{ROW\_NUMBER} number them.
  \item \textbf{(Applied)} Write the query giving each programme's average
        grade points, number of graded enrolments and number of ungraded ones,
        only for programmes with more than 5{,}000 graded enrolments. State
        for each condition which clause it belongs in, and why.
  \item \textbf{(Applied)} The university wants a merit list: the top three
        students by grade points in each programme, ties included. Write it
        with a window function, say whether you chose \sqlkw{RANK} or
        \sqlkw{DENSE\_RANK} and what the choice means for a student tied in
        third place, and estimate from the measurement how the correlated
        alternative would behave at 16{,}000 students.
\end{reviewq}
